\subsection{Canonical decomposition of a torsion module over a principal ideal domain}

Let M be a module over a commutative ring A. For each \(\alpha\in A\) let \(\mathbf{M}(\alpha)\) denote the kernel of the endomorphism \(x\mapsto\alpha x\) of M. If \(\alpha\) and \(\beta\) are two elements of A such that \(\alpha\) divides \(\beta\) , then clearly \(\mathbf{M}(\alpha)\subset\mathbf{M}(\beta)\) . In particular, as n runs through the set of rational integers \(\geqslant1\) , the submodules \(\mathbf{M}(\alpha^{n})\) form an increasing sequence; the union \(M_{\alpha}\) of the \(\mathbf{M}(\alpha^{n})\) is thus a submodule of M, consisting of those elements of M which are annihilated by some power of \(\alpha\) . For each submodule N of M, it is clear that \(N_{\alpha}=N\cap M_{\alpha}\) .

DEFINITION 1. --- Let \(\pi\) be an irreducible element of a principal ideal domain A; an A-module M is called \(\pi\) -primary if, for all \(x \in \mathbf{M}\) , there exists an integer \(n \geqslant 1\) such that \(\pi^n x = 0\) (in other words, if M is equal to the submodule \(\mathbf{M}_{\pi}\) ).

Clearly every cyclic module of the form \(\mathbf{A}/(\pi^{s})\) is \(\pi\) -primary. For an arbitrary A-module M, the submodule \(M_{\pi}\) is \(\pi\) -primary.

Lemma 1. --- Let M be a module over a principal ideal domain A; for all \(\alpha \in \mathbf{A}\) such that \(\alpha \neq 0\) , let \(\alpha = \varepsilon \sum_{i=1}^{r} \pi_i^{n(i)}\) be a decomposition of \(\alpha\) into irreducible factors (VII, p.~4). The submodule \(\mathbf{N} = \mathbf{M}(\alpha)\) of elements of M annihilated by \(\alpha\) is the direct sum of the submodules \(\mathbf{M}(\pi_i^{n(i)})\) , and the map which sends each \(x \in \mathbf{M}(\alpha)\) to its component in \(\mathbf{M}(\pi_i^{n(i)})\) has the form \(x \mapsto \gamma_i x\) ( \(\gamma_i \in \mathbf{A}\) ). Moreover

\[
\mathbf {M} \left(\pi_ {i} ^ {n (i)}\right) = \mathbf {N} \cap \mathbf {M} _ {\pi_ {i}} = \mathbf {N} _ {\pi_ {i}}.
\]

Note first that N is annihilated by \(\alpha\) , so has a natural \(\mathbf{A}/(\alpha)\) -module structure. By Prop. 4 of VII, p.~3, the canonical homomorphism from \(\mathbf{A}/(\alpha)\) into the product of the rings \(\mathbf{A}/(\pi_{i}^{n(i)})\) is a ring isomorphism; now applying Prop. 1 of VII, p.~6, we deduce that N is the direct sum of the \(\mathbf{M}(\pi_{i}^{n(i)})\) ; the projectors of this decomposition are homotheties, by VII, p.~7, Remark 2. The inclusion \(\mathbf{M}(\pi_{i}^{n(i)})\subset\mathbf{M}(\alpha)\cap\mathbf{M}_{\pi_{i}}\) is obvious; conversely, let \(x\in\mathbf{M}(\alpha)\cap\mathbf{M}_{\pi_{i}}\) . Then there exists a power \(\pi_{i}^{s}\) of \(\pi_{i}\) which annihilates x; we may assume \(s\geqslant n(i)\) ; by Bezout's identity, there exist \(\lambda,\mu\in\mathbf{A}\) such that \(\pi_{i}^{n(i)}=\lambda\pi_{i}^{s}+\mu\alpha\) , so \(\pi_{i}^{n(i)}x=0\) and so finally \(x\in\mathbf{M}(\pi_{i}^{n(i)})\) .

Lemma 2. --- Let M be a torsion module (II, p.~313) over an integral domain A. For every finite family \((x_{i})_{1 \leq i \leq n}\) of elements of M, there exists an element \(\gamma \neq 0\) in A such that the \(x_{i}\) all belong to \(\mathrm{M}(\gamma)\) .

Indeed, for each index \(i\) there exists an element \(\alpha_{i} \neq 0\) in A which annihilates \(x_{i}\) , and the element \(\gamma = \prod_{i=1}^{n} \alpha_{i}\) will fit the bill.

THEOREM 1. --- Let M be a torsion module over a principal ideal domain A; for each irreducible element \(\pi\) of A, let \(M_{\pi}\) be the submodule of M consisting of elements annihilated by some power of \(\pi\) . If P is a system of representatives of irreducible elements of A, then M is the direct sum of its submodules \(M_{\pi}\) for \(\pi \in P\) .

Every element \(x \in M\) belongs to the submodule \(\mathbf{M}(\alpha)\) for some \(\alpha \neq 0\) , so by Lemma 1 is a sum of a finite number of elements, each of which belongs to some submodule \(M_{\pi}\) . On the other hand, if \(\sum_{\pi \in P} x_{\pi} = \sum_{\pi \in P} y_{\pi}\) , where \(x_{\pi}, y_{\pi} \in M_{\pi}\) for all \(\pi \in P\) , and where all but finitely many of the \(x_{\pi}\) and \(y_{\pi}\) are zero, then Lemma 2 shows that there exists \(\gamma \neq 0\) in A such that all the \(x_{\pi}\) and \(y_{\pi}\) belong to the same submodule \(\mathbf{M}(\gamma)\) ; applying Lemma 1 to \(\mathbf{M}(\gamma)\) shows that \(x_{\pi} = y_{\pi}\) for all \(\pi \in P\) , which completes the proof.

Clearly, if \(\pi\) and \(\pi'\) are two associate irreducible elements, then \(M_{\pi} = M_{\pi'}\) ; thus, for a given module M, the submodule \(M_{\pi}\) depends only on the ideal ( \(\pi\) ) of A; it is called the \(\pi\) -primary component of the module M, and the decomposition of M as a direct sum of the \(M_{\pi}\) is called the canonical decomposition of M as a direct sum of its \(\pi\) -primary components.

COROLLARY 1. --- Every submodule N of a torsion module M is the direct sum of its submodules \(N \cap M_{\pi}\) .

This follows from the fact that \(\mathbf{N} \cap \mathbf{M}_{\pi}\) is the \(\pi\) -primary component \(\mathbf{N}_{\pi}\) of \(\mathbf{N}\) .

COROLLARY 2. --- The submodule N of the torsion A-module M is a direct factor if and only if \(N_{\pi}\) is a direct factor of \(M_{\pi}\) for every irreducible element \(\pi\) of A.

Indeed, if \(\mathbf{N}\) and \(\mathbf{N}'\) are two submodules of \(\mathbf{M}\) , then \(\mathbf{M} = \mathbf{N} \oplus \mathbf{N}'\) if and only if \(\mathbf{M}_{\pi} = \mathbf{N}_{\pi} \oplus \mathbf{N}_{\pi}'\) for every irreducible element \(\pi\) of \(\mathbf{A}\) (Cor. 1).

COROLLARY 3. --- Let N be a submodule of the torsion A-module M. If, for every irreducible element \(\pi\) of A, either \(N_{\pi} = 0\) or \((M/N)_{\pi} = 0\) , then N is a direct factor of M.

Indeed, the condition \((\mathbf{M}/\mathbf{N})_{\pi}=0\) implies \(N_{\pi}=M_{\pi}\) , and Cor. 2 applies.

An A-module M is called semi-simple if every submodule of M is a direct factor (cf.~A, VIII, § 3).

COROLLARY 4. --- Let A be a principal ideal domain which is not a field, and let M be an A-module. Then M is semi-simple if and only if M is a torsion module and \(\mathbf{M}_{\pi} = \mathbf{M}(\pi)\) for every irreducible element \(\pi\) of A.

First suppose that M is semi-simple; let \(x \in M\) and let \(\pi\) be an irreducible element of A. If N is a complement of \(A\pi x\) in M, then we can write \(x = \alpha\pi x + y\) , with \(\alpha \in A\) and \(y \in N\) ; but that implies \(y = (1 - \alpha\pi)x\) , so

\[
\pi (1 - \alpha \pi) x \in A \pi x \cap N = 0.
\]

It follows first of all that M is a torsion module; if moreover \(x \in M_{\pi}\) , then \(\pi(1 - \alpha\pi)x = 0\) , thus \(\pi x = \alpha\pi^{2}x = \alpha^{2}\pi^{3}x = \cdots = \alpha^{n}\pi^{n+1}x\) is zero and \(\mathbf{M}_{\pi} = \mathbf{M}(\pi)\) .

Conversely, by Cor. 2 it is enough to prove that an A-module M annihilated by an irreducible element \(\pi\) is semi-simple; but that is clear, since M then has a natural structure of a vector space over the field \(\mathrm{A}/(\pi)\) , and the submodules of M are precisely the vector subspaces under this structure.

Remark 1. --- Clearly the annihilator of every element \(\neq 0\) of a \(\pi\) -primary module has the form \(A\pi^k (k > 0\) an integer), since it is a principal ideal containing a power of \(\pi\) . Let \(x\) be an element of \(M\) ; for each \(\pi \in P\) , let \(x_{\pi}\) be the component of \(x\) in \(M_{\pi}\) ; the annihilator of \(x\) is the lcm of the annihilators of the nonzero \(x_{\pi}\) , but by the above it is equal in this case to the product of the annihilators of the nonzero \(x_{\pi}\) (VI, p.~16, Prop. 12 (DIV)).

PROPOSITION 2. --- If M is a finitely generated torsion module over a principal ideal domain A, then the \(\pi\) -primary components of M are zero except for a finite number of them, and the projectors of M onto these components \(M_{\pi}\) are homotheties.

This follows immediately from Lemma 1, for by Lemma 2 there exists \(\alpha \neq 0\) in A such that \(\mathbf{M} = \mathbf{M}(\alpha)\) .

Remark 2. --- By VII, p.~7, Remark 3, a finitely generated torsion A-module is cyclic if and only if each of its \(\pi\) -primary components is cyclic.

An important special case where Th. 1 and Prop. 2 apply is when A = Z; a Z-module is just an abelian group. An abelian group is called a torsion group if it is a torsion Z-module, that is if all its elements have finite order. Then P is taken to be the set of prime numbers \textgreater{} 0; for each prime number p \textgreater{} 0, an (abelian) group is said to be p-torsion if all its elements have orders which are powers of p.~In this terminology, Th. 1 shows that every torsion abelian group is a direct sum of p-torsion groups. In the case of finite groups, this also follows from I, p.~80, Th. 4.

